\section{What the local theorem contributes to a quasi-polynomial bound}
\label{sec:global}

The useful improvement is a change of parameter. If the cap graph has
$s$ non-exterior vertices, straightforward top-tree search can have
depth $s$ and a branching factor depending on the matrix size. The present
algorithm permits arbitrary $s$ at each fixed number of nonlinear
columns. It gives a polynomial-time local operator for $k=0,1,2,\ldots$
held constant, with one uniform algorithm and a computable parameter factor.

Let $N_i$ be the encoded size of the $i$th supplied graph--matrix problem,
and let $k_i$ be its residual triple-column parameter. For a sequence of
$Q$ such calls, including unsuccessful calls, the local work is bounded by
\begin{equation}
 \sum_{i=1}^{Q}2^{O(k_i\log(k_i+2))}\operatorname{poly}(N_i).
 \label{eq:cumulative-local}
\end{equation}
The polynomial exponent is independent of $k_i$. If $N_i\le n^a$ and
$k_i\le K(n)$ for every call, this is at most
\[
             Q\,n^{O(1)}2^{O(K(n)\log(K(n)+2))}.
\]
Two parameter regimes make the scale concrete:
\begin{center}
\begin{tabular}{@{}ll@{}}
\toprule
Uniform residual parameter & Local cost for polynomial-size inputs \\
\midrule
$K(n)=O(\log n/\log\log n)$ & $n^{O(1)}$ \\
$K(n)=O(\log n)$ & $n^{O(\log\log n)}$ \\
$K(n)=O((\log n)^a)$, fixed $a>0$ &
$\exp(O((\log n)^a\log\log n))\,n^{O(1)}$ \\
\bottomrule
\end{tabular}
\end{center}
Under the additional uniform deficit hypothesis of
Corollary~\ref{cor:bounded-deficit}, bounded $\Delta$ sharpens the
$K(n)=O(\log n)$ line to polynomial local time. It also permits
$K(n)=O((\log n)^2)$ within $n^{O(\log n)}$ local time.

Small values of $n$ are absorbed into constants; the first line is an
asymptotic assertion as $n$ tends to infinity.

\begin{theorem}[Conditional transfer to an entire recognition procedure]
\label{thm:global-transfer}
Suppose an exact hierarchy procedure on an $n$-crossing input satisfies:
\begin{enumerate}
\item it makes at most $\exp(O((\log n)^b))$ calls to the present operator,
including calls made after resets and on failed branches;
\item each call has encoded size at most $n^a$ and residual nonlinear
premise budget at most $O((\log n)^c)$, for fixed $a,b,c$;
\item all other work, including surface discovery, source authentication,
marked geometric transport, and final certification, has a
quasi-polynomial total bound.
\end{enumerate}
Then substituting the present exact operator preserves a
quasi-polynomial total bound. More precisely, its contribution is
\[
 \exp\!\left(O((\log n)^b+(
                   \log n)^c\log\log n+\log n)\right).
\]
\end{theorem}
\begin{proof}
Apply the general premise-budget extension of
Theorem~\ref{thm:column-fpt} to every call, and use the uniform size and
parameter bounds in \eqref{eq:cumulative-local}. Multiplication by the
number of calls adds their logarithmic costs. The sum with the stated
quasi-polynomial bound for the other operations remains quasi-polynomial.
\end{proof}

This is an accounting theorem with explicit premises. The finite solver
proves none of the three hierarchy hypotheses by itself. In particular,
a polynomial encoded-time cut does not bound the number of candidate
surfaces or the number of restored hierarchy states. A small certificate
does not imply an efficient search for it.

The rank-layer reduction provides an additional structural direction.
Under native cycle compatibility, if a scenario leaves $p$ mandatory
source groups and $q$ residual demands, then it can succeed only if
$p\le r\le p+q$. The difference $r-p$ measures the number of
additional independent initial-seed witnesses needed to justify nonlinear
activation, beyond the mandatory unary source representatives. Controlling that difference, or decomposing the exceptional
columns into independently solvable regions, may be more feasible than
bounding every raw column count.

The current source experiments deliberately retain large-parameter cases.
Stellar subdivisions give examples with increasing numbers of nonlinear
columns while the underlying manifold remains a solid torus. This makes
clear why the same topological type, or small normal weight alone, cannot
justify assuming that every labelled raw instance has small $k$.
